\chapter{前言与学习路线}

写出几个能运行的例子以后，新的问题往往接踵而来：终端能找到的库，为什么在 PyCharm 中导入失败？代码拆成多个文件以后，应当从哪个入口启动？这些问题连接着语言规则、运行环境与工程组织，也构成了本书的学习主线。

本书面向初学者，也适合希望整理零散脚本经验的读者。阅读时请把代码放回具体任务：先运行原例，再替换一组输入，对照结果变化；遇到报错，从回溯末行识别异常类型，再定位到项目中的文件和行号。

全书先从解释器与 PyCharm 入手，建立对对象、容器、控制流、函数和模块的认识；随后学习文件、数据、并发、网络及常用工具，最后通过四阶段项目练习测试、持久化与交付。每段示例都应能回答三个问题：输入是什么，规则在哪里，结果怎样验证。读完后，尝试把同样的思路用于实际任务。

\section*{开始练习前的检查}

先按第 1 章确认 Python 3、PyCharm、运行配置和断点能够正常工作，并建立 \code{.venv}。安装依赖时，使用项目选定的解释器执行 \code{python -m pip}。跟练时在 \code{requirements.txt} 中记录依赖，第 17 章再将项目元数据、依赖声明和构建配置整理到 \code{pyproject.toml}。

阅读代码前，先查阅使用说明和符号表，认清行号、交互提示符与三类提示框的用途。命名、缩进和导入顺序可对照 PEP 8；运行结果则应通过输入、输出与断点现场验证。

\par\vspace{2.3em}
\begin{center}
  {\Large\sffamily\color{PyBlue}把想法写成代码，让每一步都经得起运行。}
\end{center}
